% ===================================================================
%  اللوحة نمرة ٦ — البيت من جوّه، والعفش، والأكل، والنور.
%
%  Traduction, faite en 2026, en arabe egyptien ecrit, sur l'ido de
%  texto/io. Regles de la colonne : voir 10-lawha-01.tex. Cinq
%  pieces, cinq numerotations : les cles portent c1 a c5.
%
%  LE PLUS LONG TABLEAU DU LIVRET, ET CELUI OU L'EGYPTIEN A MEUBLE LA
%  MAISON AVEC DE L'ITALIEN ET DU TURC. Aucun de ces mots ne se dit
%  en arabe standard :
%
%    كوميدينو (28)  la table de nuit   de l'italien comodino
%    سكرتير (5)     le secretaire      de « secretaire »
%    كنبة (24)      le canape
%    برواز (30)     le cadre
%    نجفة (38)      le lustre
%    بارافان (36)   le paravent
%    فازة (3)       le vase
%    شيزلونج (47)   la chaise longue
%    بوتاجاز (19)   le fourneau        « portable gas »
%    بانيو (1)      la baignoire
%    دُش (7)         la douche
%    قيشاني (6)     le carrelage emaille
%    أباجورة (13)   l'abat-jour
%    تسريحة (1)     la coiffeuse
%
%  ET LES MOTS DE TOUS LES JOURS SE SEPARENT SANS EMPRUNTER :
%
%    عيش (12)      le pain      MSA : خبز
%    ملاية (23)    le drap      MSA : ملاءة
%    مراية (60)    le miroir    MSA : مرآة
%    معلقة (34)    la cuillere  MSA : ملعقة
%    كباية (32)    le verre     MSA : كأس
%    حنفية (48)    le robinet   MSA : صنبور
%    مقشّة (34)     le balai     MSA : مكنسة
%    طاسة (47)     la poele
%    ساطور (40)    le couperet
%    فانوس (11)    la lanterne
%    كبريت (12)    les allumettes MSA : ثقاب
%    زير (50)      la jarre a eau
%    كنكة (64)     la cafetiere
%    حلّة (25)      la marmite
%    فتلة (36)     le fil       MSA : خيط
%    ماشة (51)     les pincettes
%
%  DEUX DE CES MOTS DISENT QUELQUE CHOSE SUR L'EGYPTE ELLE-MEME. Le
%  pain s'appelle عيش, qui veut dire « la vie » : ce n'est pas une
%  image de poete, c'est le nom courant, et il n'y en a pas d'autre.
%  Et le ciel de lit (18) s'appelle ناموسية, de ناموس le moustique :
%  la planche francaise y voit un ornement, l'egyptien y voit ce a
%  quoi ca sert.
%
%  LE CIEL DE LIT ETAIT LE SEUL PIEGE DU TABLEAU. Tout le reste
%  s'enumere, et les renvois suivent l'ido sans une reprise.
% ===================================================================

%%K t06-apar-1 apar
\VUsaut{6.0mm}
\VUtitre{15.0pt}{60}{اللوحة نمرة ٦}
\VUsaut{1.4mm}
\VUtitre{11.4pt}{0}{البيت من جوّه، والعفش، والأكل،}
\VUsaut{0.4mm}
\VUtitre{11.4pt}{0}{والنور.}
\VUsaut{2.2mm}

%%K t06-apar-2 apar
البيت خلص. وعمّ يوحنا اتنقل في مسكنه الجديد، اللي إحنا عارفين ترتيبه. ودلوقت
نشوف فرشه إزاي. هنبصّ الأول على:

%%K t06-tit-1 sub
\VUpk{10.2pt}{40}{أوضة النوم.}
\VUsaut{3.0mm}

%%K t06-c1-01-1 p
1. --- إحنا جايين في وقت مش مناسب، علشان الشغّالة مشغولة بتنضّف
\VUgras{الأوضة}.

%%K t06-c1-02-1 p
2. --- على \VUgras{التسريحة} \textsuperscript{(1)} في حاجات كتير هنا وهناك،
اللي بيها الواحد بيغسل وشّه ويمشّط شعره، وبالاختصار يعمل تواليته. وهي
\VUgras{الطشت} \textsuperscript{(2)}، و\VUgras{إبريق المية} \textsuperscript{(3)}،
و\VUgras{فرشة الشعر} \textsuperscript{(4)}، و\VUgras{المشط} \textsuperscript{(5)}،
و\VUgras{الصابونة} \textsuperscript{(6)}، و\VUgras{السفنجة} \textsuperscript{(7)}، وأخيرًا
\VUgras{إزازة صغيرة} \textsuperscript{(8)} فيها معجون السنان السايل.

%%K t06-c1-03-1 p
3. --- وعلى الأرض، قدام التسريحة، أهو \VUgras{الجردل} \textsuperscript{(9)} علشان
يلمّ المية الوسخة.

%%K t06-c1-04-1 p
4. --- وإحنا شايفين في \VUgras{المدخل} \textsuperscript{(10)} كام \VUgras{شمّاعة}
\textsuperscript{(13)}، و\VUgras{شمسيتين} \textsuperscript{(11)} في \VUgras{حامل الشماسي}
\textsuperscript{(12)}.

%%K t06-c1-05-1 p
5. --- وعلى الناحية التانية من الباب في \VUgras{دولاب المراية}
\textsuperscript{(14)} اللي جوّاه \VUgras{البياضات} \textsuperscript{(15)}.

%%K t06-c1-06-1 p
6. --- خلّينا نستغلّ إن \VUgras{الشغّالة} \textsuperscript{(6)} بترتّب \VUgras{السرير}
\textsuperscript{(17)}، علشان نشوف أجزاءه المختلفة. \VUgras{هيكل السرير}
\textsuperscript{(20)} فيه \VUgras{مرتبة سوست} \textsuperscript{(21)} كويسة، وحالًا يوستينا
هتحطّ عليها \VUgras{مرتبة} \textsuperscript{(22)} صوف.

%%K t06-c1-06-1 p suite
وبعدين هتاخد من على الكراسي \VUgras{الملايتين} \textsuperscript{(23)}
و\VUgras{البطّانية} \textsuperscript{(24)}. وساعتها هتحطّ عند راس السرير
\VUgras{المسند} \textsuperscript{(25)} و\VUgras{المخدّة} \textsuperscript{(26)}، اللي هما مع
\VUgras{اللحاف الريش} \textsuperscript{(27)} على \VUgras{شرشف السرير}
\textsuperscript{(32)}. وبتستعمل ريشة علشان تشيل التراب من على \VUgras{الستاير}
\textsuperscript{(19)} ومن على \VUgras{الناموسية} \textsuperscript{(18)}، وأهو جزء من الشغل
خلص.

%%K t06-c1-07-1 p
7. --- وبعدين هيبقى لازم ترتّب شوية \VUgras{الكوميدينو} \textsuperscript{(28)}، وتلفّ
\VUgras{المنبّه} \textsuperscript{(29)}، وتحضّر \VUgras{أباجورة السرير}
\textsuperscript{(30)}،
وتشيل \VUgras{الشمعة} \textsuperscript{(31)} علشان تنضّف الشمعدان.

%%K t06-tit-2 sub
\VUsaut{5.0mm}
\VUpk{10.2pt}{40}{سبت الشغل وحاجات المدفأة.}
\VUsaut{3.0mm}

%%K t06-c1-08-1 p
8. --- و\VUgras{سبت الشغل} \textsuperscript{(34)} اللي جنب \VUgras{الكرسي الصغير}
\textsuperscript{(33)} محتاج هو كمان ترتيب. لازم تطوي \VUgras{التطريز}
\textsuperscript{(35)}، وتحطّ في \VUgras{ترابيزة الشغل} \textsuperscript{(38)}
\VUgras{الكشتبان}، و\VUgras{الفتلة} \textsuperscript{(36)}، و\VUgras{الإبر}
\textsuperscript{(37)}، و\VUgras{المقصّ} \textsuperscript{(39)}، وتقفل غطا الترابيزة
بـ\VUgras{المفتاح} \textsuperscript{(40)}.

%%K t06-c1-09-1 p
9. --- وعلى \VUgras{الترابيزة أبو رجل واحدة} \textsuperscript{(41)} اللي في النصّ،
يوستينا هتغيّر مية \VUgras{الكاراف} \textsuperscript{(42)}. وهتحطّ \VUgras{سكر} في
\VUgras{السكريّة} \textsuperscript{(43)}، وتنضّف \VUgras{ملقاط السكر} \textsuperscript{(44)}،
وتشيل \VUgras{فرشة الهدوم} \textsuperscript{(46)}.

%%K t06-c1-10-1 p
10. --- والناحية اليمين من الأوضة هتطلب منها شغل أقل، علشان
\VUgras{الشيزلونج} \textsuperscript{(47)} و\VUgras{مخدّته} \textsuperscript{(48)} في مكانهم
الصح؛ وعلشان الجو مش ساقع، مافيش حاجة اتحرّكت من حاجات \VUgras{المدفأة}
\textsuperscript{(49)}. \VUgras{الجاروف} \textsuperscript{(50)}، و\VUgras{الماشة}
\textsuperscript{(51)}، و\VUgras{حاملات الحطب} \textsuperscript{(55)}، و\VUgras{حاجز الرماد}
\textsuperscript{(56)} كلهم نضاف؛ و\VUgras{حاجز النار} \textsuperscript{(54)} ما اتحرّكش،
و\VUgras{صندوق الحطب} \textsuperscript{(52)} مليان \VUgras{حطب} \textsuperscript{(53)} كويس.

%%K t06-noto-1 noto
\VUnotes{143.2mm}{%
(*) Babo = عيّل صغير؛ بالإنجليزي \textit{baby}، وبالفرنساوي \textit{bébé}،
وبالإيطالي \textit{bambino}.%
}

%%K t06-noto-2 noto
\VUnotes{143.2mm}{%
(*) Lambrequino = بالفرنساوي \textit{lambrequin}، وبالإسباني
\textit{lambrequines}، وبالبرتغالي \textit{lambrequins}، وحتى بالألماني
والإنجليزي \textit{lambrequins}؛ يعني D. E. F. P. S.%
}

%%K t06-c1-11-1 p
11. --- ومع كده لازم تشيل التراب من على \VUgras{ساعة الحيطة} \textsuperscript{(57)}،
و\VUgras{الشمعدانات} \textsuperscript{(58)}، و\VUgras{علب الحاجات} \textsuperscript{(59)}،
وأخيرًا تمسح \VUgras{المراية} \textsuperscript{(60)}.

%%K t06-c1-12-1 p
12. --- أما \VUgras{سرير الهزّ} \textsuperscript{(61)} بتاع البيبي (*)، فهو جاهز خلاص.

%%K t06-tit-3 sub
\VUsaut{5.0mm}
\VUpk{10.2pt}{40}{الصالون.}
\VUsaut{3.0mm}

%%K t06-c2-01-1 p
1. --- ودلوقت أهو الصالون. دي \VUgras{أوضة} حلوة أوي. والمدفأة اللي في الركن
اليمين مزيّنة بـ\VUgras{تمثال صغير} \textsuperscript{(1)} رقيق، وبـ\VUgras{فازتين}
\textsuperscript{(3)} حلوين، ومكسّية بـ\VUgras{لمبريكان} \textsuperscript{(2)} قطيفة (*).

%%K t06-c2-02-1 p
2. --- وعلشان شرر النار ما يولّعش السجّادة، حطّوا قدام الموقد \VUgras{حاجز شرر}
\textsuperscript{(4)} نحاس.

%%K t06-c2-03-1 p
3. --- وجنبه، \VUgras{سكرتير} \textsuperscript{(5)} حريمي شيك مفتوح. وعليه \VUgras{ستّ
البيت} \textsuperscript{(23)} بتكتب جواباتها.

%%K t06-c2-04-1 p
4. --- والشباك متزيّن بـ\VUgras{ستاير قزاز} \textsuperscript{(6)} بـ\VUgras{أربطة}
\textsuperscript{(8)} متعلّقة في \VUgras{الخطاطيف} \textsuperscript{(7)}، وبستاير قماش تقيلة
\VUgras{مسحوبة لفوق} \textsuperscript{(9)}.

%%K t06-c2-05-1 p
5. --- وفي جوز الشباك، \VUgras{قصريّة} \textsuperscript{(11)} فيها \VUgras{نبات}
\textsuperscript{(10)} أخضر، نخلة حلوة أوي ورقها بيوصل تقريبًا لـ\VUgras{لمبة الجاز}
\textsuperscript{(12)}.

%%K t06-c2-06-1 p
6. --- و\VUgras{أباجورة} \textsuperscript{(13)} اللمبة رتّبتها ماريا،
\VUgras{الآنسة} \textsuperscript{(14)} الحلوة

%%K t06-c2-06-1 p suite
اللي على البيانو \textsuperscript{(15)}. وهي قاعدة معتدلة أوي على \VUgras{الكرسي
الدوّار} \textsuperscript{(16)}، بتذاكر قطعة موسيقى. وهي شاطرة ومرتّبة أوي، زي ما
بيبيّن \VUgras{دولاب النوتة} \textsuperscript{(17)} بتاعها اللي مرتّب تمام.

%%K t06-tit-4 sub
\VUsaut{5.0mm}
\VUpk{10.2pt}{40}{الولد الشقي.}
\VUsaut{3.0mm}

%%K t06-c2-07-1 p
7. --- وأخوها، باولوس، بيدوّر على \VUgras{كتاب} \textsuperscript{(19)} في
\VUgras{المكتبة} \textsuperscript{(18)}. وهو \VUgras{شاب} \textsuperscript{(20)} كتير الحركة
وما بيتحملش. ولمّا يتعلّق في المكتبة علشان يوصل للكتاب اللي فوق أوي، يمكن
يوقّع \VUgras{التمثال النصفي} \textsuperscript{(21)} اللي فوقها.

%%K t06-c2-08-1 p
8. --- وهو مستغلّ، علشان يعمل الغلطة دي، إن أمه مشغولة بتتكلّم مع صاحبتها،
\VUgras{ستّ} \textsuperscript{(22)} جاية تقضّي كام يوم عندها. والأم قاعدة على
\VUgras{كنبة} \textsuperscript{(24)} جنب \VUgras{الفوتيه} \textsuperscript{(25)}، وصاحبتها
على \VUgras{الكرسي} \textsuperscript{(27)} اللي مش باين منه غير \VUgras{الضهر}
\textsuperscript{(26)}.

%%K t06-c2-09-1 p
9. --- و\VUgras{البرواز} \textsuperscript{(30)} الكبير المعلّق على الحيطة فيه
\VUgras{صورة} \textsuperscript{(29)} واحدة قريبة. وفي \VUgras{الألبوم}
\textsuperscript{(32)} المحطوط على الترابيزة متجمّعة \VUgras{صور} \textsuperscript{(33)}
باقي أهل البيت. والعيال قلّبوه حالًا، علشان \VUgras{الكراسي} \textsuperscript{(31)}
لسه متجمّعة حوالين الترابيزة.

%%K t06-c2-10-1 p
10. --- و\VUgras{الفازة الصيني} \textsuperscript{(34)} البورسلين، اللي جنب
\VUgras{فتّاحة الجوابات} \textsuperscript{(35)}، اتقدّمت هدية لماريا في عيدها؛
و\VUgras{البارافان} \textsuperscript{(36)} اللي مزيّن ركن الأوضة جابه عمّها من الصين.

%%K t06-c2-11-1 p
11. --- والصالون ده، بـ\VUgras{اللوحات} \textsuperscript{(37)} الحلوة أوي المعلّقة على
\VUgras{الحاجز} \textsuperscript{(42)}، ومنوّر بـ\VUgras{النجفة} \textsuperscript{(38)} اللي
\VUgras{لمبات الكهربا} \textsuperscript{(39)} بتاعتها بتلمع بالليل

%%K t06-c2-11-1 p suite
من الفرحة، شكله مريح أوي. و\VUgras{السجّادة} \textsuperscript{(40)} الطرية المفرودة
على الباركيه، و\VUgras{الجرس الكهربا} \textsuperscript{(41)} السهل أوي، بيكمّلوا راحته.

%%K t06-tit-5 sub
\VUsaut{3.6mm}
\VUpk{10.2pt}{40}{أوضة السفرة.}
\VUsaut{3.0mm}

%%K t06-c3-01-1 p
1. --- جرس العشا دقّ. والعيلة كلها، ما عدا ماريا، متجمّعة حوالين الترابيزة.
وأوضة السفرة مدفّية كويس بـ\VUgras{الدفّاية} \textsuperscript{(1)} القيشاني الممتازة،
بـ\VUgras{ماسورتها} \textsuperscript{(2)} الرفيّعة.

%%K t06-c3-02-1 p
2. --- والأوضة دي ما فيهاش عفش كتير. وعلى طول الحيطة، معلّقة فوق
\VUgras{الكونسول} \textsuperscript{(4)}، \VUgras{نافورة صغيرة} \textsuperscript{(3)} للزينة.
وقريّب أوي، في \VUgras{البوفيه} \textsuperscript{(5)} اللي بيتحطّ عليه الأكل البارد،
وأطباق التحلاية، وحاجات السفرة المختلفة اللي مش محتاجينها حالًا.

%%K t06-c3-03-1 p
3. --- وكده إحنا شايفين على الرفّ الفوقاني

%%K t06-c3-03-1 p suite
\VUgras{فرخة محمّرة} \textsuperscript{(7)}، و\VUgras{سمكة} \textsuperscript{(8)}،
و\VUgras{طبق عالي} \textsuperscript{(10)} فيه \VUgras{فاكهة} \textsuperscript{(9)}. وتحت أهي
\VUgras{الجبنة} \textsuperscript{(11)}، و\VUgras{العيش} \textsuperscript{(12)}، و\VUgras{حامل
الزيت والخل} \textsuperscript{(13)} اللي فيه كل اللي لازم للسلطة: الزيت، والخل،
و\VUgras{الملح} \textsuperscript{(14)}، و\VUgras{الفلفل} \textsuperscript{(15)}. وأخيرًا، على
الرفّ التحتاني في \VUgras{كيكة} \textsuperscript{(17)} جنب \VUgras{المستردة}
\textsuperscript{(16)}.

%%K t06-c3-04-1 p
4. --- وساعة أوضة السفرة، اللي بيقولوا لها \VUgras{ساعة الحيطة}
\textsuperscript{(20)}، شكلها مخصوص أوي. وهي مركّبة في برواز خشب فيه كمان
\VUgras{بارومتر} \textsuperscript{(19)} و\VUgras{ترمومتر} \textsuperscript{(18)}.

%%K t06-tit-6 sub
\VUsaut{4.6mm}
\VUpk{10.2pt}{40}{خدمة العشا.}
\VUsaut{3.0mm}

%%K t06-c3-05-1 p
5. --- وكل حيطان الأوضة متلبّسة \VUgras{ألواح خشب} \textsuperscript{(21)} من سنديان
مشمّع. و\VUgras{ستارة باب} \textsuperscript{(22)} قماش بتقفل كويس الباب اللي بيودّي
للفيرندة. وهناك العيلة هتروح تشرب القهوة بعد العشا. وأهي \VUgras{الفناجين}
\textsuperscript{(24)} و\VUgras{الأطباق الصغيرة} \textsuperscript{(25)} جاهزين على
\VUgras{ترابيزة الأطباق} \textsuperscript{(23)}.

%%K t06-c3-06-1 p
6. --- وفوق الترابيزة متثبّت في السقف \VUgras{نور معلّق} \textsuperscript{(26)}،
لمبته اللمّاعة أوي بتنوّر بالليل أوضة السفرة كلها.

%%K t06-c3-07-1 p
7. --- ودلوقت خلّينا نبصّ على \VUgras{ترابيزة الأكل} \textsuperscript{(27)} نفسها.
لمّا العيلة دخلت، السفرة كانت جاهزة. وعلى \VUgras{المفرش} \textsuperscript{(28)} كان
متحطّط، في مكان كل واحد، \VUgras{طبق} \textsuperscript{(33)}، و\VUgras{فوطة}
\textsuperscript{(29)}، و\VUgras{معلقة} \textsuperscript{(34)}، و\VUgras{شوكة}
\textsuperscript{(35)}، و\VUgras{سكينة} \textsuperscript{(36)}، و\VUgras{كباية}
\textsuperscript{(32)}. وكان في كمان \VUgras{أزايز} \textsuperscript{(30)} للنبيت،
و\VUgras{كاراف} \textsuperscript{(31)} للمية.

%%K t06-c3-08-1 p
8. --- و\VUgras{الصفرجي} \textsuperscript{(39)} جاب الأول \VUgras{الشوربية}
\textsuperscript{(37)}، اللي لسه بيطلع منها بخار شوربة. ودلوقت بيقدّم \VUgras{الصنف}
\textsuperscript{(38)} الأول، صنف اللحمة. وبعد كده هيقدّم اللي سمّيناهم؛ وأخيرًا،
بعد \VUgras{صنف التحلاية}، هيستغلّ إن أسياده هيدخلوا الفيرندة علشان يشيل
حاجات السفرة ويرتّب أوضة السفرة تاني.

%%K t06-c3-08-2 p
\textit{ملحوظة.} --- \textit{الكلام الشايع اللي بيخصّ الأكل مذكور هنا بشكل
مختصر أوي. واللستة هتتكمّل في اللوحتين «الأوتيل» (نمرة ١٣) و«السوق» (نمرة
١٥).}

%%K t06-tit-7 sub
\VUsaut{4.6mm}
\VUpk{10.2pt}{40}{المطبخ.}
\VUsaut{3.0mm}

%%K t06-c4-01-1 p
1. --- وفي المطبخ، اللي بنبصّ عليه من الباب المفتوح نصّه، إحنا شايفين لخبطة
حلوة زي الصورة. وقدام \VUgras{الكانون} \textsuperscript{(1)}، اللي فيه \VUgras{النار}
\textsuperscript{(3)} بتفرقع، مركّب \VUgras{سيخ دوّار} \textsuperscript{(5)}.
و\VUgras{شوي} \textsuperscript{(7)} حلو \VUgras{متسيّخ} \textsuperscript{(6)} وبينضج.

%%K t06-c4-02-1 p
2. --- و\VUgras{المنفاخ} \textsuperscript{(8)} و\VUgras{جاروف الفحم}
\textsuperscript{(9)}، اللي اتستعملوا حالًا، فضلوا على الأرض زيهم زي
\VUgras{الحطب} \textsuperscript{(10)}.

%%K t06-c4-03-1 p
3. --- وفوق المدفأة مرتّب؛ \VUgras{الفانوس} \textsuperscript{(11)}، و\VUgras{علبة
الكبريت} \textsuperscript{(13)} اللي فيها \VUgras{الكبريت} \textsuperscript{(12)}،
و\VUgras{الشمعدان} \textsuperscript{(14)}، و\VUgras{حامل الشمعة} \textsuperscript{(15)}
بـ\VUgras{الشمعة} \textsuperscript{(16)}، كلهم في مكانهم. بس \VUgras{جردل الفحم}
\textsuperscript{(18)} الكبير، المليان \VUgras{فحم} \textsuperscript{(17)} اللي
\VUgras{الطبّاخة} \textsuperscript{(23)} حالًا ملته من البدروم، فضل في نصّ المطبخ.

%%K t06-c4-04-1 p
4. --- وبالحقّ، الستّ المسكينة لازم تستعجل علشان الغدا يبقى جاهز في وقته.
ودلوقت هي جنب \VUgras{البوتاجاز} \textsuperscript{(19)} وبتقلّب \VUgras{صلصة}
\textsuperscript{(24)} لازم ما تتحرقش.

%%K t06-c4-05-1 p
5. --- وفي \VUgras{الفرن} \textsuperscript{(20)} بتتخبز كيكة حضّرتها لأكل العيال.
وفوق البوتاجاز معلّقة \VUgras{المغرفة} \textsuperscript{(21)} و\VUgras{المصفاية}
\textsuperscript{(22)}.

%%K t06-tit-8 sub
\VUsaut{4.0mm}
\VUpk{10.2pt}{40}{الحلل.}
\VUsaut{3.0mm}

%%K t06-c4-06-1 p
6. --- و\VUgras{الحلل} \textsuperscript{(25)}، المعلّقة في آخر الأوضة، نضيفة أوي.
و\VUgras{الكسرولّات} \textsuperscript{(26)} النحاس الحلوة، و\VUgras{إبريق اللبن}
\textsuperscript{(27)}، و\VUgras{المصفاية الصغيرة} \textsuperscript{(28)}،
و\VUgras{المبشرة} \textsuperscript{(29)} اتنضّفوا بعناية.

%%K t06-c4-07-1 p
7. --- و\VUgras{علبة الملح} \textsuperscript{(30)} محطوطة على دولاب الأطباق جنب
\VUgras{براد الشاي} \textsuperscript{(31)}، و\VUgras{صفيحة الزيت} \textsuperscript{(32)}.
و\VUgras{الدرج} \textsuperscript{(33)} على الأغلب فيه المعالق والسكاكين بتاعة المطبخ.

%%K t06-c4-08-1 p
8. --- و\VUgras{المقشّة} \textsuperscript{(34)} و\VUgras{الريشة} \textsuperscript{(35)} في
ركن. والطبّاخة استعملت حالًا \VUgras{القرمة} \textsuperscript{(36)}، وسابتها جنب الشباك.

%%K t06-c4-09-1 p
9. --- و\VUgras{فوطة الإيد} \textsuperscript{(37)} معلّقة على الحيطة، جنب
\VUgras{القمع} \textsuperscript{(38)}. وفي الناحية دي من المطبخ متحطّطة
\VUgras{الشوّاية} \textsuperscript{(39)}، و\VUgras{الساطور} \textsuperscript{(40)}،
و\VUgras{لوح التقطيع} \textsuperscript{(41)}. وبعد كده، فوق الساطور، على \VUgras{رفّ}
\textsuperscript{(42)}، في \VUgras{قدور} \textsuperscript{(43)}، و\VUgras{حلّة السلق}
\textsuperscript{(44)} بـ\VUgras{غطاها} \textsuperscript{(45)}، و\VUgras{الحلّة الكبيرة}
\textsuperscript{(46)}. و\VUgras{الطاسة} \textsuperscript{(47)} معلّقة جنبهم.

%%K t06-c4-10-1 p
10. --- ومافيش حاجة بتلمع في الركن ده غير \VUgras{الحنفية} \textsuperscript{(48)}
النحاس. و\VUgras{الحوض} \textsuperscript{(49)}، اللي عليه \VUgras{الزير}
\textsuperscript{(50)} التخين، على الأغلب مريح أوي بدولابه الصغير، اللي بيتحفظ فيه
\VUgras{برميل الخل} \textsuperscript{(51)} بـ\VUgras{حنفيته} \textsuperscript{(52)}،
و\VUgras{صفيحة الزبالة} \textsuperscript{(53)}، و\VUgras{الأطباق الوسخة}
\textsuperscript{(54)}.

%%K t06-tit-9 sub
\VUsaut{4.0mm}
\VUpk{10.2pt}{40}{حاجات تانية موجودة في المطبخ.}
\VUsaut{3.0mm}

%%K t06-c4-11-1 p
11. --- \VUgras{الطشت الكبير} \textsuperscript{(55)} اللي بيتغسل فيه
\VUgras{الخضار} \textsuperscript{(58)}، و\VUgras{الحلّة} \textsuperscript{(56)} الحديد
المحطوطة على \VUgras{البلاط} \textsuperscript{(57)}، على الأغلب ليهم كمان مكان في
الدولاب ده.

%%K t06-c4-12-1 p
12. --- والطبّاخة أكيد مش ناقصها بوتاجازات، علشان أهو كمان، على ركن
الترابيزة، \VUgras{بوتاجاز غاز} \textsuperscript{(59)} عليه مية بتسخن في كسرولّة
صغيرة. و\VUgras{البخار} \textsuperscript{(60)} اللي طالع منها بيقول لنا إنها على
الأغلب بتغلي. ويمكن المية دي تتستعمل للقهوة، علشان أهي على الطرف التاني من
الترابيزة \VUgras{الكنكة} \textsuperscript{(64)} و\VUgras{مطحنة البن}
\textsuperscript{(63)}.

%%K t06-c4-13-1 p
13. --- والبوتاجاز موصول بـ\VUgras{شعلة الغاز} \textsuperscript{(62)}
بـ\VUgras{خرطوم كاوتش} \textsuperscript{(61)} طويل.

%%K t06-c4-14-1 p
14. --- و\VUgras{الشغّالة باليومية} \textsuperscript{(66)} اللي جاية تساعد الطبّاخة،
بتحضّر الخضار للعشا.

%%K t06-c4-14-1 p suite
ولمّا يبقوا اتقشّروا واتغسلوا، هتحطّهم في \VUgras{الطشت} \textsuperscript{(65)}
لحدّ ما ييجي وقت الطبيخ.

%%K t06-tit-10 sub
\VUsaut{4.0mm}
{\centering\fontsize{10.2pt}{10.2pt}\selectfont\textls[40]{\scshape الحمّام.} \textit{(أوضة الاستحمام.)}\par}
\VUsaut{3.0mm}

%%K t06-c5-01-1 p
1. --- ما بقاش فاضل لنا غير قلّة أدب واحدة علشان نعرف أوض الشقة الرئيسية:
إننا نشوف تركيب الحمّام. وده مش هياخد وقت طويل، علشان هو مش كبير، وحالًا
هنعرف إن \VUgras{البانيو} \textsuperscript{(1)} عنده \VUgras{سخّان} \textsuperscript{(2)}
كويس، وإن \VUgras{صحن الصابون} \textsuperscript{(3)} إيد اللي بيستحمّى بتوصل له، وإن
\VUgras{حامل الفوط} \textsuperscript{(5)} أكيد هيسهّل تنشيف \VUgras{الفوط}
\textsuperscript{(4)}.

%%K t06-c5-02-1 p
2. --- والأوضة دي حيطانها متلبّسة \VUgras{قيشاني} \textsuperscript{(6)}، اللي سمح
بتركيب \VUgras{دُش} \textsuperscript{(7)} من غير ما نخاف إن المية، اللي بترشّ وهو
شغّال، تبوّظ الحيطان. و\VUgras{طشت} \textsuperscript{(8)} زنك صغير، بيستعمل لغسل
الرجلين، بيكمّل العفش.

\VUsaut{6.0mm}
\VUfilet{22mm}
